\section{Experimental Setup}
\label{sec:experiments}

\subsection{Datasets}
% TODO: brief RSNA + SPIDER description.
\textbf{RSNA 2024 Lumbar Spine Degenerative Classification.}
$\sim$2{,}000 patients, sagittal T2 stack per IVD, 3 conditions
(spinal canal, left/right foraminal) $\times$ 3 severity classes
(Normal/Moderate/Severe). Patient-level 80/20 split, $\textsc{seed}=42$,
$1{,}942$ validation IVDs.

\textbf{SPIDER.} $\sim$1{,}400 IVDs from a separate hospital cohort, 8
disease labels (Pfirrmann grade, Modic, Disc narrowing/herniation/bulging,
Spondylolisthesis, UP/LOW endplate). Used both as a fully held-out
zero-shot test set and as a supervised transfer-learning target.

\subsection{Evaluation Protocol}
% Mean F1 macro + Mean Accuracy primary; domain-specific (Severe F1,
% per-tier zero-shot F1) reported alongside.
We adopt \textbf{Mean F1 macro} (per-condition macro-F1 averaged across
conditions) and \textbf{Mean Accuracy} as primary metrics in all
cross-phase comparisons. Domain-specific metrics — Severe F1 on RSNA
and per-tier zero-shot F1 (easy/medium/hard) on SPIDER — are reported
alongside but never as the sole metric.

\subsection{Phases}
\begin{itemize}
  \item \textbf{Phase 1+2} (RSNA supervised): Baseline vs CBAM vs
        Hybrid (CBAM + frozen BiomedCLIP). Class imbalance handled
        with focal loss + sqrt-frequency class weights.
  \item \textbf{Phase 3} (SPIDER zero-shot): No SPIDER training. The
        Hybrid model trained on RSNA encodes 8 SPIDER prompts at
        inference and predicts. Compared against ``Naked BiomedCLIP''
        (same architecture without the RSNA-trained backbone or
        projection MLP) to isolate the effect of the 3D specialist
        branch.
  \item \textbf{Phase 4} (SPIDER supervised transfer): 4-way
        comparison — Vanilla, CBAM, Hybrid frozen, Hybrid full
        fine-tune — on 4 SPIDER conditions (Pfirrmann 5-cls,
        Modic 4-cls, Disc Narrowing 2-cls, Spondylolisthesis 2-cls).
\end{itemize}
